\section{Conclusion}

We could say something about: (But not sure whether this makes sense.)
- central identity of classical IIT and relation to cause-effect repertoire
- possibility to consider simpler cause-effect repertoires
- point to process theories again and to further examples (quantum / leaking) in other paper
- say something about meaning of the crucial steps (probably rather not)
- Repertoires: In IIT 3.0, this map is defined in terms of dynamical properties of the system.
- q-shape describing the \emph{quality} of the systems experience.		
- List new mathematical problems in our Outlook (existence, non-uniqueness, behaviour over coarse-graining, stability)

Outlook / usefulness of this reconstruction:
- shown: solve math problems
- could include lower bound thing
- control over the behaviour of the theory
- our major hope: guidance in the further development of the theory


For later use:
- Product of proto-spaces = compositional structure of experience (?)
- Proto space first, map is second.
In classical IIT, this is intertwined